\section{Model mismatch (exploratory)}
\label{sec:mismatch}

\subsection{Design}

The results of Section~\ref{sec:heldout} were obtained on deterministic objectives.
Under those conditions the planner's internal model is exact, so the plan formed at
the initial state is already the optimal prediction and replanning cannot acquire
new information. Two readings of $\gate{C3} \approx 0$ are therefore possible:

\begin{description}
  \item[Reading 1] Feedback has no value to begin with.
  \item[Reading 2] This task family is predictable, so feedback was not needed.
\end{description}

The pivotal question is not ``is the model nonlinear?'' but \emph{``is the future
state unpredictable at the time the initial plan is formed?''} We therefore changed
only the samples the optimizer observes, for \emph{the same model and data}:

\begin{center}
\begin{tabular}{ll}
\toprule
Regime & Description \\
\midrule
R1 full-batch & gradients and HVPs over the full dataset; deterministic \\
R2 minibatch  & a fixed-seed batch sequence; the sample changes every step \\
\bottomrule
\end{tabular}
\end{center}

R2 places subsampled curvature inside Newton--CG, a setting known to be difficult
in its own right \citep{byrd2011stochastic}. Our aim is not to resolve that
difficulty but to see whether the value of replanning changes when the planner's
internal predictive model is wrong.

The problem is a two-layer MLP ($4{,}869$ parameters) on a five-class
classification task with $512$ samples generated by a fixed teacher network. The
evaluation objective is always the full dataset, so $\JE$ remains comparable across
regimes.

Batches advance \emph{only after an actual step}. They do not advance during the
planner's look-ahead simulation, so the planner predicts the future using the
current batch. Letting it see future batches would amount to oracle knowledge of
the data and would defeat the purpose of R2.

\Needspace{6\baselineskip}
\subsection{Results}

\begin{table}[!ht]
  \centering
  % 이 파일은 scripts/make_tables.py 가 생성한다. **손으로 수정하지 않는다.**
% 표: micro-neural absolute
% raw: headroom_micro-neural_step_size_fixed_b8_0bec1125.jsonl
% raw 의 SHA-256 은 paper/evidence_map.md 에 있다.
% 공개 저장소에는 raw 가 없다. results/public/*.csv 에서 --public-dir 로 만들면
% 이 파일과 바이트 단위로 같아진다 (scripts/verify_public_results.py).
\begin{tabular}{lrrr}
\toprule
Controller & R1 full-batch & R2 batch 128 & R2 batch 64 \\
\midrule
\texttt{best\_static} & 5.220 & 3.765 & 3.029 \\
\texttt{best\_open\_loop} & 3.299 & 2.100 & 2.984 \\
\texttt{heuristic} & 5.105 & 2.162 & 0.986 \\
\texttt{onestep\_narrow} & 19.886 & 4.296 & 3.468 \\
\texttt{onestep\_absolute} & 19.460 & 4.642 & 3.732 \\
\texttt{committed\_Q4\_narrow} & 20.491 & -0.402 & 1.086 \\
\texttt{shrinking\_Q4\_narrow} & 20.395 & 3.185 & 2.757 \\
\bottomrule
\end{tabular}

  \caption{Absolute median $\JE$ on the micro-neural problem, $n = 3$ per regime.
  Exploratory.}
  \label{tab:micro-absolute}
\end{table}

\begin{table}[!ht]
  \centering
  % 이 파일은 scripts/make_tables.py 가 생성한다. **손으로 수정하지 않는다.**
% 표: micro-neural paired delta (exploratory, n=3 per regime)
% raw: headroom_micro-neural_step_size_fixed_b8_0bec1125.jsonl
% raw 의 SHA-256 은 paper/evidence_map.md 에 있다.
% 공개 저장소에는 raw 가 없다. results/public/*.csv 에서 --public-dir 로 만들면
% 이 파일과 바이트 단위로 같아진다 (scripts/verify_public_results.py).
\begin{tabular}{lrrr}
\toprule
Comparison & R1 full-batch & R2 batch 128 & R2 batch 64 \\
\midrule
shrinking $-$ best\_static & $+15.176$ (3/3) & $-0.900$ (1/3) & $-0.277$ (1/3) \\
shrinking $-$ onestep & $+0.547$ (2/3) & $-1.111$ (1/3) & $-0.716$ (1/3) \\
shrinking $-$ committed & $-0.095$ (1/3) & $+2.941$ (3/3) & $+1.666$ (3/3) \\
committed $-$ onestep & $+0.470$ (2/3) & $-4.444$ (0/3) & $-2.322$ (0/3) \\
\bottomrule
\end{tabular}

  \caption{Paired differences by regime. \emph{We do not quote confidence
  intervals or $p$-values} at $n = 3$; the parenthesized figure is the count of
  positive pairs.}
  \label{tab:micro-paired}
\end{table}

\subsection{Why $\gate{C3} > 0$ in R2}
\label{sec:c3-r2}

The absolute values distinguish the two readings:

\begin{center}
\begin{tabular}{lrrr}
\toprule
Controller & R1 & R2 (batch 64) & R2 (batch 128) \\
\midrule
\ctrl{committed} & 20.491 & 1.086 & $-0.402$ \\
\ctrl{shrinking} & 20.395 & 2.757 & 3.185 \\
\bottomrule
\end{tabular}
\end{center}

In R2, \ctrl{committed} falls below even \ctrl{best\_static} and trails
\ctrl{onestep} by $2.3$ to $4.4$~nat. The rejection rate is decisive.

\begin{figure}[!ht]
  \centering
  \includegraphics[width=\textwidth]{figure3_model_mismatch.png}
  \caption{Exploratory micro-neural results under model mismatch, $n = 3$ per
  regime. Panel (a) compares terminal improvement across the full-batch and
  minibatch regimes. Panel (b) gives the rejection rate of \ctrl{committed}:
  $0.00$ in R1, $0.79$ at batch 128, $0.66$ at batch 64, so a large fraction of
  the plan formed at the initial state is rejected by later batches.}
  \label{fig:mismatch}
\end{figure}

That is, $\gate{C3} > 0$ does not mean \ctrl{shrinking} became better; it means
\emph{holding a stale plan is costly}.

At the same time, in R2 the planner falls below both the tuned constant and the
inexpensive one-step controller. State-conditioned control does help in R2
(\ctrl{onestep} $3.468$ versus \ctrl{best\_static} $3.029$), but one-step greedy
control already provides it at $1/146$ of the planner's search cost (batch-64
regime medians: \ctrl{onestep} $1{,}879$~GE, \ctrl{shrinking} $275{,}286$~GE).

\subsection{Limitations of this experiment}

\begin{itemize}
  \item $n = 3$ per regime. We do not quote CIs or $p$-values.
  \item The transition across three batch sizes is not monotone: batch 128 scores
        below batch 64.
  \item Cross-regime absolute comparisons are not FLOP-normalized
        (Section~\ref{sec:ge}).
  \item One model, one dataset.
\end{itemize}

The R1 value $\gate{C2} = +0.547$ agrees in sign with the $+0.456$ of
Section~\ref{sec:heldout} ($n = 40$, $p < 0.0001$), but becomes $-0.080$ under the
alternative acceptance criterion. The $n = 40$ result is the stronger one and
$n = 3$ does not overturn it. \emph{We report the disagreement itself as a result.}
